%% codeact-loop.tex -- the CodeAct agent loop of the first journal paper,
%% redrawn with figkit. Colour codes the role: green for what the user supplies,
%% blue for what the environment offers, pink for what the model produces,
%% grey for what the runtime executes.
%%
%% The cards are placed first, then the group boxes fitted around them, and only
%% then the single blocks, each pinned with \fkrow to the group or card it talks
%% to, so User, Inputs and LLM share one line and Python Code and the interpreter
%% share another. Every arrow between them is exactly horizontal by construction.
%%
%% Requires figkit.sty. Input inside a figure* environment:
%%   \input{figures/codeact-loop}
\providecommand{\fkfigwidth}{\linewidth}
\resizebox{\fkfigwidth}{!}{%
\begin{tikzpicture}[font=]

% ===================================================================
% Cards
% ===================================================================
% --- what the user supplies ---------------------------------------
\node[fkcardtitlec=appleGreen] (qTitle) at (3.6,12.0) {Question};
\node[fkprose, text width=4.6cm] (qBody) at (3.6,11.55)
      {e.g. `What is the total volume of concrete used for the foundation?'};
\node[fkcardtitlec=appleGreen] (bTitle) at (3.6,9.4) {BIM Model};
\node[fkprose, text width=4.6cm] (bBody) at (3.6,8.95)
      {e.g. \texttt{duplex\_house/arc.ifc}};

% --- what the environment offers ----------------------------------
\node[fkcardtitlec=appleBlue] (tTitle) at (3.6,5.6) {Tools};
\node[fkprose, text width=4.6cm] (tBody) at (3.6,5.15)
      {e.g. \texttt{categorize\_element\_by\_properties}};
\node[fkcardtitlec=appleBlue] (sTitle) at (3.6,3.2) {Special tool};
\node[fkprose, text width=4.6cm] (sBody) at (3.6,2.75)
      {\texttt{query\_ifcopenshell\_documentation}};

% --- what the model produces --------------------------------------
\node[fkcardtitlec=applePink] (thTitle) at (14.8,6.9) {Thoughts};
\node[fkprose, text width=4.4cm] (thBody) at (14.8,6.45)
      {e.g. `In order to find the total volume of concrete used, I need to [...]'};
\node[fkcardtitlec=applePink] (pcTitle) at (14.8,4.3) {Python Code};
\node[fkprose, text width=4.4cm] (pcBody) at (14.8,3.85)
      {e.g. \texttt{"import ifcopenshell model = ifcopenshell.open("...}};

% --- what comes back out ------------------------------------------
\node[fkcardtitlec=applePink] (aTitle) at (13.6,15.0) {Answer};
\node[fkprose, text width=4.4cm] (aBody) at (13.6,14.55)
      {e.g. `33\,m$^3$ of concrete were used for the foundation'};

\begin{scope}[on background layer]
  \node[fkcardc=appleGreen, fit=(qTitle)(qBody)]   (cardQ)  {};
  \node[fkcardc=appleGreen, fit=(bTitle)(bBody)]   (cardB)  {};
  \node[fkcardc=appleBlue,  fit=(tTitle)(tBody)]   (cardT)  {};
  \node[fkcardc=appleBlue,  fit=(sTitle)(sBody)]   (cardS)  {};
  \node[fkcardc=applePink,  fit=(thTitle)(thBody)] (cardTh) {};
  \node[fkcardc=applePink,  fit=(pcTitle)(pcBody)] (cardPc) {};
  \node[fkcardc=applePink,  fit=(aTitle)(aBody)]   (cardA)  {};
\end{scope}

% ===================================================================
% Group boxes around the cards
% ===================================================================
\begin{pgfonlayer}{grouping}
  \node[fkgroup={appleGreen}{Inputs},      fit=(cardQ)(cardB)]   (grpIn)  {};
  \node[fkgroup={appleBlue}{Resources},    fit=(cardT)(cardS)]   (grpRes) {};
  \node[fkgroup={applePink}{Final output}, fit=(cardA)]          (grpAns) {};
\end{pgfonlayer}

% ===================================================================
% Blocks, each on the line of what it exchanges with
% ===================================================================
\node[fkblock={appleGray}{2.2cm}]     (user)  at (\fkrow{0}{grpIn.west})    {User};
\node[fkblock={applePink}{3.4cm}]     (llm)   at (\fkrow{14.2}{grpIn.east}) {LLM};
\node[fkblock={appleDarkGray}{3.6cm}] (pyint) at (\fkrow{23.2}{cardPc.east})
      {Python\\Interpreter};
\node[fkblock={appleDarkGray}{3.6cm}] (evres) at (\fkcol{pyint.north}{8.6})
      {Evaluation\\Result};

% the loop box is drawn before the Outputs group it contains, so its fill does
% not cover the inner box: on one layer, the later node wins
\begin{pgfonlayer}{grouping}
  \node[fkgroup={appleDark}{CodeAct loop}, inner sep=0.85cm,
        fit=(llm)(cardTh)(cardPc)(pyint)(evres)] (loop) {};
  \node[fkgroup={applePink}{Outputs}, fit=(cardTh)(cardPc)] (grpOut) {};
\end{pgfonlayer}

% ===================================================================
% Wiring
% ===================================================================
% the user asks, and the question plus the model reach the LLM
\draw[fkarrowc=appleGreen] (user.east)  -- (grpIn.west);
\draw[fkarrowc=appleGreen] (grpIn.east) -- (llm.west);
% the environment offers tools to the model and a runtime to execute in
\draw[fkarrowc=appleBlue, fkelbow] (grpRes.east) -- ++(1.9,0)
      |- ($(llm.west)+(0,-0.40)$);
\draw[fkarrowc=appleBlue, fkelbow] ($(grpRes.east)+(1.9,0)$) |- ($(loop.south)+(0,-0.7)$)
      -| (pyint.south);
% the model emits a thought and the code that follows from it
\draw[fkarrowc=applePink] (llm.south) -- (llm.south |- grpOut.north);
% the code is executed and its result goes back to the model
\draw[fkarrowc=appleDarkGray] (cardPc.east) -- (pyint.west);
\draw[fkarrowc=appleDarkGray] (pyint.north) -- (evres.south);
\draw[fkarrowc=appleDarkGray, fkelbow] (evres.west) -- ++(-1.1,0) |- (llm.east);
% once the model is done it answers, and the answer reaches the user
\draw[fkarrowc=applePink] (llm.north) -- (llm.north |- grpAns.south);
\draw[fkarrowc=applePink, fkelbow] (grpAns.west) -| (user.north);

\end{tikzpicture}}
